%----------------------------------------------------------------------------------------
%	PACKAGES AND OTHER DOCUMENT CONFIGURATIONS
%----------------------------------------------------------------------------------------

\documentclass[10pt,stdletter,dateno,sigleft]{newlfm} % Extra options: 'sigleft' for a left-aligned signature, 'stdletternofrom' to remove the from address, 'letterpaper' for US letter paper - consult the newlfm class manual for more options

%\usepackage{charter} % Use the Charter font for the document text
\usepackage{graphics}

\usepackage{standalone}

%% with this patch, the newlfm should work (BEGIN)
\usepackage{etoolbox}
\makeatletter
\patchcmd{\@zfancyhead}{\fancy@reset}{\f@nch@reset}{}{}
\patchcmd{\@set@em@up}{\f@ncyolh}{\f@nch@olh}{}{}
\patchcmd{\@set@em@up}{\f@ncyolh}{\f@nch@olh}{}{}
\patchcmd{\@set@em@up}{\f@ncyorh}{\f@nch@orh}{}{}
\makeatother
%% with this patch, the newlfm should work (END)


\newsavebox{\Luiuc}\sbox{\Luiuc}{\parbox[b]{1.75in}{\vspace{0.5in}
\includegraphics[width=1.3\linewidth]{/Users/hectorbahamonde/Job_Market/logo.png}}} % Company/institution logo at the top left of the page
\makeletterhead{Uiuc}{\Lheader{\usebox{\Luiuc}}}

\newlfmP{sigsize=50pt} % Slightly decrease the height of the signature field
\newlfmP{addrfromphone} % Print a phone number under the sender's address
\newlfmP{addrfromemail} % Print an email address under the sender's address
\PhrPhone{\texttt{p}} % Customize the "Telephone" text
\PhrEmail{\texttt{e}} % Customize the "E-mail" text

\lthUiuc % Print the company/institution logo

% email input begin
\newread\fid
\newcommand{\readfile}[1]% #1 = filename
{\bgroup
  \endlinechar=-1
  \openin\fid=#1
  \read\fid to\filetext
  \loop\ifx\empty\filetext\relax% skip over comments
    \read\fid to\filetext
  \repeat
  \closein\fid
  \global\let\filetext=\filetext
\egroup}
\readfile{/Users/hectorbahamonde/Bibliografia_PoliSci/email.txt}
% email input end

\usepackage{color} % for colors

\usepackage{hyperref}
\hypersetup{
    %bookmarks=true,         % show bookmarks bar?
    unicode=false,          % non-Latin characters in Acrobat’s bookmarks
    pdftoolbar=true,        % show Acrobat’s toolbar?
    pdfmenubar=true,        % show Acrobat’s menu?
    pdffitwindow=true,     % window fit to page when opened
    pdfstartview={FitH},    % fits the width of the page to the window
    pdftitle={My title},    % title
    pdfauthor={Author},     % author
    pdfsubject={Subject},   % subject of the document
    pdfcreator={Creator},   % creator of the document
    pdfproducer={Producer}, % producer of the document
    pdfkeywords={keyword1} {key2} {key3}, % list of keywords
    pdfnewwindow=true,      % links in new window
    colorlinks=true,       % false: boxed links; true: colored links
    linkcolor=blue,          % color of internal links (change box color with linkbordercolor)
    citecolor=blue,        % color of links to bibliography
    filecolor=blue,      % color of file links
    urlcolor=blue           % color of external links
}

\hypersetup{
  colorlinks = true,
  urlcolor = blue,
  pdfpagemode = UseNone
}


%%% bib begin
\usepackage[american]{babel}
\usepackage{csquotes}
%\usepackage[style=chicago-authordate,doi=false,isbn=false,url=false,eprint=false]{biblatex}

\usepackage[authordate,isbn=false,doi=false,url=false,eprint=false]{biblatex-chicago}
\DeclareFieldFormat[article]{title}{\mkbibquote{#1}} % make article titles in quotes
\DeclareFieldFormat[thesis]{title}{\mkbibemph{#1}} % make theses italics

\AtEveryBibitem{\clearfield{month}}
\AtEveryCitekey{\clearfield{month}}

\addbibresource{/Users/hectorbahamonde/Bibliografia_PoliSci/library.bib} 
\addbibresource{/Users/hectorbahamonde/Bibliografia_PoliSci/Bahamonde_BibTex2013.bib} 

% USAGES
%% use \textcite to cite normal
%% \parencite to cite in parentheses
%% \footcite to cite in footnote
%% the default can be modified in autocite=FOO, footnote, for ex. 
%%% bib end

%----------------------------------------------------------------------------------------
%	YOUR NAME AND CONTACT INFORMATION
%----------------------------------------------------------------------------------------

%\namefrom{\vspace{-3cm}Hector Bahamonde, PhD} % Name

\addrfrom{
{\vspace{-2cm}}\\ % Date
[12pt]
{\bf Hector Bahamonde, PhD}\\
{\color{blue}{\bf Senior Researcher}}\\
\emph{{\color{blue}INVEST}}\\
{\color{blue}University of Turku}\\
Turku, Finland
}

\phonefrom{+358 (0) 40-362-1438} % Phone number

\emailfrom{\href{mailto:hibano@utu.fi}{hibano@utu.fi} \\ 
\texttt{w}: \href{http://www.hectorbahamonde.com}{www.HectorBahamonde.com}\\
\today}% Link to last version


%----------------------------------------------------------------------------------------
%	ADDRESSEE AND GREETING/CLOSING
%----------------------------------------------------------------------------------------

%\greetto{\vspace{-1cm}Dear Members of the Search Committee,} % Greeting text
%\closeline{\vspace{-1cm}\includegraphics[width=3cm]{/Users/hectorbahamonde/Administracion/signature.pdf} \hspace{-4cm}} % Closing text




%----------------------------------------------------------------------------------------

\begin{document}
\begin{newlfm}

%----------------------------------------------------------------------------------------
%	LETTER CONTENT
%----------------------------------------------------------------------------------------

%\vspace{-1cm}
\vspace{-0.3cm}I am writing to apply for the tenure-track position available at the Faculty of Social Sciences of the University of Turku. Currently, I am a Senior Researcher at the INVEST flagship unit, University of Turku. In this cover letter I will outline my motivations to the task and vision for the future development of research, teaching and societal interaction in the discipline.

First, I'll briefly describe my trajectory. Before joining University of Turku in 2021, I was an Assistant Professor (tenure-track) for three years at Universidad de O'Higgins (UOH), Chile. Previously, I had served one year as a post-doctoral fellow at Tulane University (New Orleans, Louisiana, U.S.). I received my PhD in Political Science from Rutgers University (New Brunswick, New Jersey, U.S.) in the spring of 2017. {\bf My research and teaching interests are the political economy of inequality and democratic development, as well as quantitative/statistical, experimental and quasi-experimental methods}. 




% bio and my research in general
%I obtained my B.A. in Political Science at Catholic University of Chile. Then I received a full fellowship from Rutgers University-New Brunswick, NJ, United States, to pursue my PhD in Political Science, which I earned in May 2017. I spent one year as a post-doctoral fellow at Tulane University-New Orleans, LA, U.S. {\bf My research explores the political economy of inequality and democratic development. I am particularly passionate and versed in quantitative/statistical, experimental and quasi-experimental methods}. 

% Experience at INVEST (general)
I've been lucky enough to be a Senior Researcher at INVEST since September 2021. This March 2022 will mark 7 months of hard-work which come with a profound sense of intellectual stimulation and professional accomplishment. In this period of time, I've published ``\href{https://link.springer.com/article/10.1057/s41269-020-00174-4}{Still for Sale: The Micro-Dynamics of Vote Selling in the United States, Evidence from a List Experiment}'' (A1) and ``\href{https://www.sciencedirect.com/science/article/pii/S0176268021000471?via%3Dihub}{Inclusive Institutions, Unequal Outcomes: Democracy, State Capacity, and Income Inequality}'' (A1). But just as important, I've enjoyed so much the process of collaborating with different researchers from INVEST, the Turku School of Economics and the Political Science Unit at the University of Turku. Not only that, I am actively collaborating with scholars from the \AA bo Akademi (Turku) who share similar research interests with me---I explain these collaboration instances further below. In sum, I have embraced with joy this intellectual community which has always been welcoming, warm and positively challenging to me. As I explain below, the TT position would allow me to continue contributing even more to the University of Turku.

% Strengthen 
This position would allow me to strengthen even more the methodological and substantive capabilities of the Faculty of Social Sciences and the INVEST ecosystem. In particular, the post will allow me to contribute with research and teaching outputs centered around experimental and quasi-experimental designs, as well as, quantitative/statistical methods. This contribution will not only materialize in peer-reviewed papers, but also in the classroom, with service, by obtaining external funds, and by joining the community outside the walls of our campus (e.g. societal interaction). Let me elaborate some more. 

% Research: Experimental and quasi-experimental
Experimental and quasi-experimental social science is very much what I have always done, and continue doing. While randomization is the gold standard to establish causal relationships, quasi-experimental designs often times offer good-enough second-bests options. For instance, using matching techniques I published in 2018 ``\href{https://journals.sub.uni-hamburg.de/giga/jpla/article/view/1121/1128}{Aiming Right at You: Group versus Individual Clientelistic Targeting in Brazil},'' (A1) where I argued that vote-buying was not only a function of earnings but also income inequality. However, every time that's possible---and ethically feasible---experiments should be conducted. Following this reasoning, we developed an economic experiment in ``\href{https://github.com/hbahamonde/Economic_Experiment_Vote_Selling/raw/master/abstract.txt}{Electoral Risk and Vote Buying, Introducing Prospect Theory in the Experimental Study of Clientelism}.'' In this paper we find that vote buying is most likely when parties are probable \emph{winners} (not \emph{losers}, as traditional research explains). {\bf The paper is currently under review in \emph{Electoral Studies}}. Taking these findings as a starting point, Prof. Janne Tukiainen (Turku School of Economics) and myself are currently collaborating in a new design that considers income inequality and political risk. Partnering with the \href{https://pcrclab.utu.fi/?page_id=894&lang=en}{PCRClab} will be fundamental to accomplish this goal. 

% Research: Statistical
Quantitative/statistical work is another area that I also actively exploiting (and has always been). For instance, after publishing the ``Inclusive Institutions, Unequal Outcomes'' paper, where we use time-series methods, Heikki 




% develop in one line these parenthesis.
% classroom (both undergraduate and graduate)
% service (supervising students and coordinating seminars as Research Director and demonstrating leadership by organizing conferences/workshops, editing journals)
% obtaining external funds
% outreach (international cooperation, societal interaction)
top-notch research

experimental research and teaching



with my teaching and research experience, both substantively and methodologically. 

% Family
Since my wife and two small children are German, my main goal is to {\bf begin an academic career in Europe}. I am looking forward to start from the bottom of the academic hierarchy. Also I have immediate availability. In this cover letter I will highlight some details about my research and teaching agendas.



{More information, \href{http://www.hectorbahamonde.com/teaching/}{\texttt{syllabi}}, my \href{http://github.com/hbahamonde/Job_Market/raw/master/Bahamonde_Research_Statement.pdf}{\texttt{research}}, \href{http://github.com/hbahamonde/Job_Market/raw/master/Bahamonde_Teaching_Statement.pdf}{\texttt{teaching}} and \href{http://github.com/hbahamonde/Job_Market/raw/master/Bahamonde_Diversity_Statement.pdf}{\texttt{diversity}} statements, as well as other \href{http://www.hectorbahamonde.com/research/}{\texttt{papers}} are available on my website: \href{http://www.hectorbahamonde.com}{\texttt{www.HectorBahamonde.com}}. Thank you for considering my application. I look forward to hearing from you.\unskip}


\vspace{-10cm}


%----------------------------------------------------------------------------------------

\end{newlfm}
\end{document}
